\documentclass{beamer}
\usetheme{Madrid}
\usepackage[utf8]{inputenc}
\usepackage[vietnamese]{babel}
\usepackage{booktabs}
\usepackage{amsmath}

% \usepackage[backend=biber,style=authoryear]{biblatex}
% \addbibresource{references.bib}

\title[ReSCORE]{%
\textsc{ReSCORE}: Label-free Iterative Retriever Training for Multi-hop QA
}
\subtitle{\small Lee et al, ACL 2025.~\cite{lee2025rescore}}

\author{Nhóm 4}
\date{\today}

\begin{document}

\begin{frame}
    \titlepage
\end{frame}

\section{MHQA và Iterative RAG}
\begin{frame}{Mục lục}
    \tableofcontents[currentsection]
\end{frame}

\begin{frame}{Multi-Hop Question Answering (MHQA)}

MHQA là bài toán trả lời câu hỏi yêu cầu mô hình suy luận qua nhiều bước và kết hợp thông tin từ nhiều tài liệu.

\begin{block}{Question}
Tác giả của \textit{Hannibal and Scipio} học ở đâu?
\end{block}

\begin{itemize}
    \item Tìm tác giả của \textit{Hannibal and Scipio}
    \[
    \textit{Hannibal and Scipio} \rightarrow \text{Elizabeth Carter}
    \]

    \item Tìm nơi Elizabeth Carter được giáo dục
    \[
    \text{Elizabeth Carter} \rightarrow \text{educated at home by her father}
    \]
\end{itemize}

\textbf{Answer} : At home, by her father.

\end{frame}

\begin{frame}{Iterative RAG}
    \begin{figure}
        \centering
        \includegraphics[width=0.4\linewidth]{image1.png}
        \label{fig:placeholder}
    \end{figure}
\end{frame}

\begin{frame}{Related Work}
\small

\begin{itemize}
    \item \textbf{Training Retriever for RAG}:
    Các phương pháp như REALM~\cite{guu2020realm}, DPR~\cite{karpukhin2020dpr},
    RAG~\cite{lewis2020rag} và Atlas~\cite{izacard2023atlas}
    tập trung train retriever để lấy tài liệu tốt hơn cho QA/RAG.

    \item \textbf{Iterative RAG}:
    IRCoT~\cite{trivedi2023ircot}, FLARE~\cite{jiang2023flare},
    Self-RAG~\cite{asai2023selfrag} và Adaptive-RAG~\cite{jeong2024adaptive}
    lặp lại quá trình retrieve--generate để xử lý câu hỏi nhiều bước.

    \item \textbf{Training with LLM Supervision}:
    Promptagator~\cite{dai2023promptagator}, InPars~\cite{bonifacio2022inpars}
    và Syntriever~\cite{kim2025syntriever}
    dùng LLM để tạo dữ liệu hoặc nhãn giả cho retriever.
\end{itemize}

\begin{block}{Thiếu sót}
Các hướng trên chưa giải quyết trực tiếp việc train dense retriever cho
\textbf{iterative multi-hop QA} khi query thay đổi qua từng iteration và không có document labels thủ công.
\end{block}

\end{frame}

% \begin{frame}{Related Work: Training Retriever for RAG}

% \begin{itemize}
%     \item \textbf{REALM}~\cite{guu2020realm}: kết hợp retriever vào quá trình pre-training language model.
    
%     \item \textbf{DPR}~\cite{karpukhin2020dpr}: train dense retriever bằng cặp question--positive passage cho open-domain QA.
    
%     \item \textbf{RAG}~\cite{lewis2020rag}: kết hợp dense retriever với generator để sinh answer dựa trên retrieved passages.
    
%     \item \textbf{Atlas}~\cite{izacard2023atlas}: cải thiện few-shot QA bằng retriever + language model quy mô lớn.
% \end{itemize}

% \begin{alertblock}{Hạn chế}
% Các phương pháp này chủ yếu train retriever cho single-hop hoặc dựa trên nhãn query--document có sẵn; chưa xử lý tốt cho iterative MHQA.
% \end{alertblock}

% \end{frame}

% \begin{frame}{Related Work: Iterative RAG}

% \begin{itemize}
%     \item \textbf{IRCoT}~\cite{trivedi2023ircot}: xen kẽ retrieval với chain-of-thought reasoning cho câu hỏi nhiều bước.
    
%     \item \textbf{FLARE}~\cite{jiang2023flare}: chủ động retrieve thêm thông tin khi mô hình dự đoán thiếu chắc chắn.
    
%     \item \textbf{Self-RAG}~\cite{asai2023selfrag}: mô hình tự quyết định khi nào cần retrieve và tự đánh giá generation.
    
%     \item \textbf{Adaptive-RAG}~\cite{jeong2024adaptive}: điều chỉnh chiến lược retrieval theo độ phức tạp của câu hỏi.
% \end{itemize}

% \begin{alertblock}{Hạn chế}
% Các phương pháp này  không tập trung vào việc train dense retriever cho phù hợp với target domain.
% \end{alertblock}

% \end{frame}

% \begin{frame}{Related Work: Training with LLM Supervision}

% \begin{itemize}
%     \item \textbf{Promptagator}~\cite{dai2023promptagator}: dùng LLM sinh synthetic queries để train retriever cho domain mới.
    
%     \item \textbf{InPars}~\cite{bonifacio2022inpars}: tạo query giả từ documents bằng LLM để cải thiện neural retriever.
    
%     \item \textbf{Syntriever}~\cite{kim2025syntriever}: train retriever bằng synthetic data sinh từ black-box LLM.
    
%     \item \textbf{Self-RAG}~\cite{asai2023selfrag}: dùng reflection/special tokens để học khi nào cần retrieve và đánh giá output.
% \end{itemize}

% \begin{block}{Điểm chung}
% Dùng LLM để sinh dữ liệu để train/inference.
% \end{block}

% \end{frame}

% \begin{frame}{Các phương pháp Retriever hiện tại}

% % \textbfG{Giải pháp trong MHQA hiện tại}

% \vspace{0.2cm}

% \textbf{Sparse Retriever (BM25 - Robertson, et al.~\cite{robertson2009bm25})}
% \begin{itemize}
% \item Dựa trên \textbf{keyword matching}
% \item Không cần huấn luyện
% \item Dễ áp dụng cho MHQA
% \item Nhưng giới hạn khả năng biểu diễn ngữ nghĩa
% \end{itemize}

% \vspace{0.2cm}

% \textbf{Dense Retriever (Contriever - Izacard, et al.~\cite{izacard2022contriever} )}

% \begin{itemize}
% \item Sử dụng \textbf{query embedding} và \textbf{document embedding}
% \item Hiệu quả hơn BM25 trong nhiều bài toán retrieval
% \item Nhưng cần dữ liệu gán nhãn để fine-tune
% \end{itemize}

% \vspace{0.2cm}

% \end{frame}

% \begin{frame}{Thách thức khi train dense retriever}
% \begin{itemize}
%     \item Mỗi iteration có thể dùng một query khác nhau.
%     \item Query được reformulate bởi LLM, nên thay đổi theo từng model.
%     \item Vì vậy, rất khó gán nhãn thủ công:
%     \[
%     q^{(i)} \rightarrow d_j^{(i)} \text{ relevant hay không?}
%     \]
% \end{itemize}

% -> việc tạo nhãn relevance cho document ở từng bước là rất tốn kém.

% \end{frame}

\section{Giải pháp ReSCORE}
\begin{frame}{Mục lục}
    \tableofcontents[currentsection]
\end{frame}

\begin{frame}{Giải pháp: ReSCORE}
\textbf{Retriever Supervision with Consistency and Relevance} :

\begin{itemize}
    \item Tạo \textbf{pseudo-GT} để train retriever.
\end{itemize}

    Ngoài ra còn đảm bảo được :
    
\begin{itemize}
    \item Document - \textbf{relevance(liên quan)} đến câu hỏi.
    \item Document - \textbf{consistency(nhất quán)} với đáp án đúng.
\end{itemize}

\end{frame}

% \begin{frame}{Iterative RAG Framework}

% \textbf{Bài toán MHQA}

% Cho:

% \[
% q:\text{Question}
% \]

% \[
% D=\{d_1,d_2,\dots,d_n\}
% \]

% là tập tài liệu.

% Mục tiêu:

% \[
% D^*=\{d_1^*,...,d_h^*\}\subseteq D
% \]

% Trong đó:

% \begin{itemize}
% \item $D^*$: tập tài liệu cần thiết để trả lời
% \item Hệ thống phải tìm đủ $D^*$ để sinh đáp án
% \end{itemize}

% \end{frame}

% \begin{frame}{Iterative RAG Framework}

% Tại vòng lặp thứ $i$:

% \[
% q^{(i)}
% \rightarrow
% D^{(i)}
% =
% \{d_1^{(i)},...,d_k^{(i)}\}
% \]

% Retriever lấy $k$ document liên quan.

% \vspace{0.2cm}

% LLM sẽ:

% \begin{itemize}
% \item Nếu đủ thông tin:

% \[
% D^{(i)}
% \rightarrow
% Answer
% \]

% \item Nếu chưa đủ:

% \[
% D^{(i)}
% \rightarrow
% t^{(i)}
% \rightarrow
% q^{(i+1)}
% \]

% Sinh \textit{Thought} để tóm tắt thông tin và tạo query mới.
% \end{itemize}

% Quá trình lặp cho tới khi tìm đủ thông tin để sinh đáp án.

% \end{frame}

\begin{frame}{Train Retriever for Iteractive RAG} 

% \begin{block}{Như đã đề cập}
% Vì không có sẵn labels cho relevant documents vậy nên ta cần 1 phương pháp để giải quyết.
% \end{block}
Lấy  phân phối sau làm \textbf{Preudo - Ground Truth} :
\[
Q^{(i)}_{\mathrm{LM}}\left(d^{(i)}_j \mid q^{(i)}\right)
\]
Mang ý nghĩa là Likelihood retrieve document \(d^{(i)}_j\) cho bởi query \(q^{(i)}\) tại iteration \(i\) do Language Model sinh ra.

\end{frame}

\begin{frame}{Formally}

Với câu hỏi q, a là câu trả lời tương ứng cho bởi document $d_j^{(i)}$..

Biểu diễn của phân phối :
\[
Q_{\mathrm{LM}}^{(i)}\left(d_j^{(i)} \mid q\right)
\propto
P_{\mathrm{LM}}^{(i)}\left(a, q \mid d_j^{(i)}\right)
\]

\[
=
P_{\mathrm{LM}}^{(i)}\left(q \mid d_j^{(i)}\right)
\cdot
P_{\mathrm{LM}}^{(i)}\left(a \mid q, d_j^{(i)}\right)
\]

Trong đó, \(P_{\mathrm{LM}}\) là xác suất của một chuỗi token sinh ra từ LLM.

\begin{itemize}
    \item \(P_{\mathrm{LM}}(q \mid d)\)
    \(\Rightarrow\) đo \textbf{relevance}.

    \item \(P_{\mathrm{LM}}(a \mid q,d)\)
    \(\Rightarrow\) đo \textbf{consistency}.
\end{itemize}

\end{frame}

\begin{frame}{Training Loss: KL Divergence}
\begin{block}{Mục tiêu}
Tối thiểu hóa KL divergence giữa pseudo-GT distribution và retrieval distribution.
\end{block}

\[
\sum_{n=1}^{N}
\sum_{i=0}^{\eta_n}
D_{\mathrm{KL}}
\left(
Q_{\mathrm{LM}}^{(i)}
\left(\mathcal{D}^{(i)} \mid q_n^{(i)}\right)
\;\|\;
P_R^{(i)}
\left(\mathcal{D}^{(i)} \mid q_n^{(i)}\right)
\right)
\]

\begin{itemize}
    \item \(N\): số lượng câu hỏi trong tập train.
    \item \(\eta_n\): số iteration của câu hỏi thứ \(n\).
    \item \(P_R^{(i)}\) là phân phối xác suất của các document được retriever tính tại iteration \(i\)..
\end{itemize}

\end{frame}

\begin{frame}{Phân phối retrieval \(P_R^{(i)}\)}
Phân phối \(P_R^{(i)}\) được tính bằng Softmax trên độ tương đồng giữa
query embedding và document embedding.

\[
P_R^{(i)}\left(d_j^{(i)} \mid q_n^{(i)}\right)
=
\mathrm{Softmax}
\left(
\bm{d}_j^{(i)} \cdot \bm{q}_n^{(i)}
\right)
\]

\[
\bm{d}_j^{(i)} = \mathrm{Embed}_{doc}\left(d_j^{(i)}\right),
\qquad
\bm{q}_n^{(i)} = \mathrm{Embed}_{query}\left(q_n^{(i)}\right)
\]

\begin{itemize}
    \item \(\bm{d}_j^{(i)}\): vector embedding của document \(d_j^{(i)}\).
    \item \(\bm{q}_n^{(i)}\): vector embedding của query \(q_n^{(i)}\).
\end{itemize}

\end{frame}

\begin{frame}{\(P_{\mathrm{LM}}\)}

\begin{figure}
    \centering
    \includegraphics[width=0.8\linewidth]{ảnh_PLM.png}
    \caption{Vai trò của các thành phần trên MuSiQue dataset}
    \label{fig:placeholder}
\end{figure}

\end{frame}

% \begin{frame}{Traning trong thực tế}

% \begin{itemize}
%     \item Vì tính \(Q_{\mathrm{LM}}^{(i)}\) trên toàn bộ database \(\mathcal{D}\) là quá tốn kém:
%     \[
%     |\mathcal{D}| \text{ rất lớn}
%     \quad \text{nên} \quad
%     \text{chi phí tính toán của LLM rất lớn}
%     \]
%     \item Vì vậy, tại mỗi iteration \(i\), chỉ lấy top-\(M\) documents theo score của retriever:
%     \[
%     M \ll |\mathcal{D}|
%     \]
%     \item Sau đó chỉ tính \(Q_{\mathrm{LM}}^{(i)}\) trên tập top-\(M\) này.
% \end{itemize}

% \end{frame}

\begin{frame}{Overview of ReSCORE}
\begin{figure}
    \centering
    \includegraphics[width=1\linewidth]{Ovv_ReSCORE.png}
    \caption{ReSCORE}
    \label{fig:placeholder}
\end{figure}
\end{frame}

\begin{frame}{Ưu điểm của ReSCORE}
\small

\begin{itemize}
    \item \textbf{Không cần document labels thủ công}: dùng LLM để tạo pseudo-GT cho retriever.
    
    \item \textbf{Train dense retriever hiệu quả hơn}: retriever học ưu tiên documents vừa liên quan đến question, vừa hỗ trợ answer đúng.
    
    \item \textbf{Có thể cải thiện nhiều iterative RAG frameworks}: thực hiện finetune cho nhiều method dense retrieve.
\end{itemize}

\end{frame}

% \begin{frame}{Nhược điểm của ReSCORE}
% \small

% \begin{itemize}
%     \item \textbf{Chi phí tính toán cao}: cần LLM để chấm nhiều documents ở nhiều iterations.
    
%     \item \textbf{Phụ thuộc vào chất lượng LLM}: nếu LLM đánh giá sai relevance/consistency, pseudo-GT sẽ nhiễu.
    
%     \item \textbf{Out-of-Distribution generalization challenge}: hiệu quả kém nhiều khi đưa ra dự đoán với các dataset khác tính chất.
    
% \end{itemize}

% \end{frame}

\section{Thực nghiệm}
\begin{frame}{Mục lục}
    \tableofcontents[currentsection]
\end{frame}


\begin{frame}{Bộ dữ liệu Musique}
\small

\begin{block}{Giới thiệu}
Musique (\textbf{Multihop Questions via Question Composition}) là bộ dữ liệu Question Answering tiếng Anh được thiết kế để đánh giá khả năng suy luận đa bước (Multi-hop Reasoning).
\end{block}

\begin{center}
\begin{tabular}{p{0.32\linewidth} p{0.50\linewidth}}
\toprule
\textbf{Đặc điểm} & \textbf{Mô tả} \\
\midrule
Ngôn ngữ & Tiếng Anh \\
Loại bài toán & Multi-hop Question Answering \\
Nguồn dữ liệu & Wikipedia \\
Đáp án & Short Answer \\
Thách thức chính & Retrieval + Multi-hop Reasoning \\
\bottomrule
\end{tabular}
\end{center}

\vspace{0.2cm}

\end{frame}

\begin{frame}{Kết quả chạy thực nghiệm với data của paper}
\begin{table}[h]
\centering
\Large
\renewcommand{\arraystretch}{1.25}
\setlength{\tabcolsep}{22pt}

\resizebox{\textwidth}{!}{
\begin{tabular}{llccc}
\toprule[1.5pt]
\textbf{Source} & \textbf{Dataset} & \textbf{EM} & \textbf{F1} \\
\midrule
Paper  & MuSiQue & 23.4 & 32.7 \\
Ours  & MuSiQue & \textbf{24.6} & \textbf{33.5} \\
\bottomrule[1.5pt]
\end{tabular}
}

\caption{So sánh kết quả thực nghiệm trên tập dữ liệu MuSiQue}
\label{tab:musique_compare}
\end{table}
\end{frame}

%============================================================
% \begin{frame}{Kết quả trên Musique}
% \small

% \begin{center}
% \begin{tabular}{lcc}
% \toprule
% \textbf{Model} & \textbf{EM} & \textbf{F1} \\
% \midrule

% Adaptive-RAG & 23.6 & 31.8 \\
% IRCoT & 22.0 & 31.8 \\
% Adaptive-Note & 13.2 & 24.2 \\
% Our Baseline (Llama-3.1-8B+BM25) & 15.2 & 23.6 \\
% Our Baseline (Llama-3.1-8B+Contriever) & 15.2 & 23.8 \\

% \midrule
% \textbf{Proposed Method} &
% \textbf{28.23} &
% \textbf{39.55} \\
% \bottomrule
% \end{tabular}
% \end{center}

% \vspace{0.3cm}

% \end{frame}
% %===========================================
% \begin{frame}{Nhận xét}
%     \item Phương pháp đề xuất đạt kết quả tốt nhất trên bộ dữ liệu Musique với \textbf{EM = 28.23} và \textbf{F1 = 39.55}.
    
%     \item So với Adaptive-RAG, EM tăng từ \textbf{23.6} lên \textbf{28.23} (+19.5\%) và F1 tăng từ \textbf{31.8} lên \textbf{39.55} (+24.4\%).

%     \item Kết quả cho thấy hệ thống không chỉ cải thiện khả năng truy xuất tài liệu liên quan mà còn nâng cao hiệu quả suy luận đa bước (multi-hop reasoning).

%     \item Mức cải thiện đặc biệt rõ rệt trên Musique, một bộ dữ liệu đòi hỏi kết hợp thông tin từ nhiều tài liệu để đưa ra đáp án chính xác.
% \end{frame}
%=========================================== 

\begin{frame}{Bộ Dữ Liệu mới: ViMQA}

\begin{itemize}
    \item ViMQA là bộ dữ liệu hỏi đáp tiếng Việt được giới thiệu tại LREC 2022.
    
    \item Bao gồm hơn 10,000 cặp câu hỏi -- trả lời được xây dựng từ Wikipedia tiếng Việt.

    \item Các câu hỏi yêu cầu suy luận đa bước (Multi-hop Reasoning) trên nhiều đoạn văn.

    \item Cung cấp Supporting Facts để giải thích đáp án.

\end{itemize}

\end{frame}

%==================================================


\begin{frame}{Chuẩn Hóa Câu Trả Lời (Normalization)}

\begin{itemize}
    \item Một số dự đoán đúng về mặt ngữ nghĩa nhưng khác ngôn ngữ với Ground Truth.
    
    \item Ví dụ:
    \begin{itemize}
        \item Europe $\rightarrow$ Châu Âu
        \item China $\rightarrow$ Trung Quốc
        \item Yes $\rightarrow$ Đúng
        \item No $\rightarrow$ Không
    \end{itemize}

    \item Thực hiện bước chuẩn hóa để ánh xạ các đáp án tiếng Anh sang tiếng Việt trước khi đánh giá.

    \item Giúp giảm lỗi định dạng và phản ánh chính xác hơn chất lượng mô hình.
\end{itemize}

\end{frame}
%==================================================
\begin{frame}{So Sánh Kết Quả Trên Tập Test}

\begin{table}[h]
\centering
\small
\begin{tabular}{lcc}
\hline
\textbf{Phương pháp} & \textbf{EM} & \textbf{F1} \\
\hline
XLM-RoBERTa$_{Base}$      & 43.76 & 59.38 \\
XLM-RoBERTa$_{Large}$     & 49.75 & 65.64 \\
InfoXLM$_{Base}$          & 49.05 & 65.76 \\
InfoXLM$_{Large}$         & 49.75 & 65.29 \\
\hline
Rescore (Raw)             & 30.23 & 45.49 \\
Rescore (Normalized)      & \textbf{49.94} & \textbf{66.75} \\
\hline
\end{tabular}
\caption{Kết quả trên tập test của ViMQA}
\end{table}

\vspace{0.3cm}

\begin{itemize}
    \item Chuẩn hóa câu trả lời giúp tăng:
    \begin{itemize}
        \item EM: +19.21 điểm
        \item F1: +21.26 điểm
    \end{itemize}
    \item F1 sau chuẩn hóa vượt các mô hình baseline được báo cáo.
\end{itemize}

\end{frame}
%==================================================
\begin{frame}{So sánh độ khó của Musique và ViMQA}
% \small

% \begin{block}{Phương pháp}
% Đánh giá độ khó thông qua log suy luận của cùng một hệ thống QA:
% \begin{itemize}
%     \item Số bước suy luận trung bình (Reasoning Steps)
%     \item Số tài liệu được sử dụng trung bình (Retrieved Documents)
%     \item Tỷ lệ câu hỏi cần nhiều hơn một bước suy luận
% \end{itemize}
% \end{block}

\begin{center}
\begin{tabular}{lcc}
\toprule
\textbf{Metric} & \textbf{Musique} & \textbf{ViMQA} \\
\midrule
Avg. Reasoning Steps & \textbf{1.731} & 1.167 \\

Avg. Retrieved Documents & \textbf{1.748} & 1.167 \\

Multi-thought Rate & \textbf{18.97\%} & 4.85\% \\
\bottomrule
\end{tabular}
\end{center}

\vspace{0.2cm}

\begin{itemize}
    \item Musique cần nhiều hơn \textbf{1.48$\times$} bước suy luận.
    \item Musique cần nhiều hơn \textbf{1.50$\times$} tài liệu hỗ trợ.
    \item Tỷ lệ trả lời câu hỏi multi-hop question hơn gần \textbf{4 lần}.
\end{itemize}

\end{frame}

% \begin{frame}{Kết luận}
% Các chỉ số đều cho thấy Musique đòi hỏi nhiều suy luận hơn ViMQA:
% \begin{itemize}
%     \item Nhiều bước reasoning hơn.
%     \item Cần truy xuất nhiều tài liệu hơn.
%     \item Tỷ lệ câu hỏi multi-hop cao hơn đáng kể.
% \end{itemize}


% \end{frame}

% \begin{frame}{\(P_{\mathrm{LM}}\)}

% \begin{figure}
%     \centering
%     \includegraphics[width=0.8\linewidth]{ảnh_PLM.png}
%     \caption{Vai trò của các thành phần trên MuSiQue dataset}
%     \label{fig:placeholder}
% \end{figure}

% \end{frame}

\begin{frame}{Nhược điểm của ReSCORE}
\small

\begin{itemize}
    \item \textbf{Chi phí tính toán cao}: cần LLM để chấm nhiều documents ở nhiều iterations.
    
    \item \textbf{Phụ thuộc vào chất lượng LLM}: nếu LLM đánh giá sai relevance/consistency, pseudo-GT sẽ nhiễu.
    
    \item \textbf{Out-of-Distribution generalization challenge}: hiệu quả kém nhiều khi đưa ra dự đoán với các dataset khác tính chất.
    
\end{itemize}

\end{frame}

\begin{frame}{Demo}
    \begin{figure}
        \centering
        \includegraphics[width=1\linewidth]{ảnh.png}
        \caption{1 minh họa cho Demo}
        \label{fig:placeholder}
    \end{figure}
\end{frame}

\begin{frame}[allowframebreaks]{References}
\footnotesize
\begin{thebibliography}{9}

\bibitem{lee2025rescore}
Dosung Lee, et al.
\newblock \textit{ReSCORE: Label-free Iterative Retriever Training for Multi-hop Question Answering with Relevance-Consistency Supervision}.
\newblock ACL 2025.

\bibitem{robertson2009bm25}
Stephen Robertson and Hugo Zaragoza.
\newblock \textit{The Probabilistic Relevance Framework: BM25 and Beyond}.
\newblock Foundations and Trends in Information Retrieval, 2009.

\bibitem{izacard2022contriever}
Gautier Izacard, Mathilde Caron, Lucas Hosseini, Sebastian Riedel, Piotr Bojanowski, Armand Joulin, and Edouard Grave.
\newblock \textit{Unsupervised Dense Information Retrieval with Contrastive Learning}.
\newblock Transactions on Machine Learning Research, 2022.

\end{thebibliography}
\end{frame}

\begin{frame}[allowframebreaks]{References}
\footnotesize
\begin{thebibliography}{99}

\bibitem{guu2020realm}
Kelvin Guu, Kenton Lee, Zora Tung, Panupong Pasupat, and Ming-Wei Chang.
\newblock \textit{REALM: Retrieval-Augmented Language Model Pre-Training}.
\newblock ICML, 2020.

\bibitem{karpukhin2020dpr}
Vladimir Karpukhin, Barlas Oguz, Sewon Min, Patrick Lewis, Ledell Wu, Sergey Edunov, Danqi Chen, and Wen-tau Yih.
\newblock \textit{Dense Passage Retrieval for Open-Domain Question Answering}.
\newblock EMNLP, 2020.

\bibitem{lewis2020rag}
Patrick Lewis, Ethan Perez, Aleksandra Piktus, Fabio Petroni, Vladimir Karpukhin, Naman Goyal, Heinrich K{\"u}ttler, Mike Lewis, Wen-tau Yih, Tim Rockt{\"a}schel, Sebastian Riedel, and Douwe Kiela.
\newblock \textit{Retrieval-Augmented Generation for Knowledge-Intensive NLP Tasks}.
\newblock NeurIPS, 2020.

\bibitem{izacard2023atlas}
Gautier Izacard, Patrick Lewis, Maria Lomeli, Lucas Hosseini, Fabio Petroni, Timo Schick, Jane Dwivedi-Yu, Armand Joulin, Sebastian Riedel, and Edouard Grave.
\newblock \textit{Atlas: Few-shot Learning with Retrieval Augmented Language Models}.
\newblock JMLR, 2023.

\end{thebibliography}
\end{frame}

\begin{frame}[allowframebreaks]{References}
\footnotesize
\begin{thebibliography}{99}

\bibitem{trivedi2023ircot}
Harsh Trivedi, Niranjan Balasubramanian, Tushar Khot, and Ashish Sabharwal.
\newblock \textit{Interleaving Retrieval with Chain-of-Thought Reasoning for Knowledge-Intensive Multi-Step Questions}.
\newblock ACL, 2023.

\bibitem{jiang2023flare}
Zhengbao Jiang, Frank F. Xu, Luyu Gao, Zhiqing Sun, Qian Liu, Jane Dwivedi-Yu, Yiming Yang, Jamie Callan, and Graham Neubig.
\newblock \textit{Active Retrieval Augmented Generation}.
\newblock EMNLP, 2023.

\bibitem{asai2023selfrag}
Akari Asai, Zeqiu Wu, Yizhong Wang, Avirup Sil, and Hannaneh Hajishirzi.
\newblock \textit{Self-RAG: Learning to Retrieve, Generate, and Critique through Self-Reflection}.
\newblock ICLR, 2024.

\bibitem{jeong2024adaptive}
Soyeong Jeong, Jinheon Baek, Sukmin Cho, Sung Ju Hwang, and Jong C. Park.
\newblock \textit{Adaptive-RAG: Learning to Adapt Retrieval-Augmented Large Language Models through Question Complexity}.
\newblock NAACL, 2024.

\end{thebibliography}
\end{frame}

\begin{frame}[allowframebreaks]{References}
\footnotesize
\begin{thebibliography}{99}

\bibitem{dai2023promptagator}
Zhuyun Dai, Vincent Y. Zhao, Ji Ma, Yi Luan, Jianmo Ni, Jing Lu, Anton Bakalov, Kelby Hall, and Ming-Wei Chang.
\newblock \textit{Promptagator: Few-shot Dense Retrieval from 8 Examples}.
\newblock ICLR, 2023.

\bibitem{bonifacio2022inpars}
Luiz Henrique Bonifacio, Vitor Jeronymo, Hugo Queiroz Abonizio, Israel Campiotti, Marzieh Fadaee, Roberto Lotufo, and Rodrigo Nogueira.
\newblock \textit{InPars: Data Augmentation for Information Retrieval using Large Language Models}.
\newblock arXiv, 2022.

\bibitem{kim2025syntriever}
Minchan Kim, et al.
\newblock \textit{Syntriever: How to Train Your Retriever with Synthetic Data from Black-box LLMs}.
\newblock Findings of NAACL, 2025.

\end{thebibliography}
\end{frame}

\end{document}
